\frontslide{Hachage extensible}


%------------------------------------------------------------
\begin{frame}{Hachage extensible}

Le hachage extensible permet de réorganiser la table de hachage en fonction
des insertions et suppressions. 

\vfill

Idée de base\,: la fonction de hachage $h$ est fixe, \alert{mais} 
on utilise les $n$ premiers bits du résultat $h(c)$ 
pour s'adapter à la taille de la collection.

\vfill

Par rapport à la version de base du hachage, on ajoute deux contraintes.

\begin{redbullet}
  \item le nombre d'entrées dans le répertoire est une puissance
          de 2
  \item la fonction $h$ donne toujours un entier sur 4 octets (32 bits)  
\end{redbullet}

\vfill
\begin{blocgauche}
La taille du répertoire tend à croître rapidement et peut poser problème.
\end{blocgauche}
\end{frame}
%------------------------------------------------------------
\begin{frame}{Exemple\,: le hachage des 16 films}

On ne montre que le premier octet.

\vfill

\begin{tabular}{l|c}
titre & $h$(titre)   \\
\hline
Vertigo     & 01110010  \\
Brazil      & 10100101  \\
Twin Peaks  & 11001011  \\
Underground & 01001001  \\
Easy Rider  & 00100110 \\
Psychose    &  01110011 \\
Greystoke   &  10111001 \\
Shining     &  11010011 \\
\hline
\end{tabular}
\vfill

La fonction $h$ est immuable. 

\end{frame}
%------------------------------------------------------------
\begin{frame}{Construction de la table}

On départ on utilise seulement le premier bit
de la fonction
\begin{redbullet}
  \item Deux valeurs possibles\,: 0 et 1
  \item Donc deux entrées, et deux fragments
  \item L'affectation d'un enregistrement dépend
        du premier bit de sa fonction de hachage  
\end{redbullet}

\vfill

  \figSlideWithSize{hachext2}{8cm}

Avec les 5 premiers films.
\end{frame}

\begin{frame}{Insertions}

Supposons 3 films par fragment. L'insertion de Psychose
(valeur 01110011) entraîne le débordement du premier fragment.

\vfill

\begin{redbullet}
  \item On double la taille du répertoire: quatre entrées 00, 01, 10 et 11.
  \item On alloue un nouveau fragment pour l'entrée 01
  \item Les entrées 10 et 11 pointent sur le \alert{même} fragment.
\end{redbullet}

\vfill

\begin{blocgauche}
Le répertoire grandit, mais dans l'espace de stockage on ajoute seulement
un nouveau fragment.
\end{blocgauche}


\end{frame}
%------------------------------------------------------------
\begin{frame}{Illustration}

Les enregistrements anciennement stockés dans le fragment 0 sont répartis
dans les fragments 00 et 01 en fonction de leur second bit.
\vfill

  \figSlideWithSize{hachext3}{9cm}

\vfill

Réorganisation \alert{locale}, donc coût acceptable.

\end{frame}
%------------------------------------------------------------
\begin{frame}{Insertions suivantes}

Deux cas.

\vfill

\alert{Cas 1}\,: on insère dans un fragment plein, mais plusieurs
           entrées pointent dessus
  
 $\Rightarrow$ on alloue un nouveau fragment, et on répartit les pointeurs

\vfill

\alert{Cas 2}\,:  on insère dans un fragment plein, associé à une seule entrée
 
     
 $\Rightarrow$ on double à nouveau le nombre d'entrées

\vfill

\begin{blocgauche}
Simple, mais la croissance du  répertoire est un potentiel problème.
\end{blocgauche}

\end{frame}
%------------------------------------------------------------
\begin{frame}{Greystoke 10111001 et Shining 11010011}

Le fragment 1 déborde. Il était référencé par deux entrées du répertoire.

\vfill

  \figSlideWithSize{hachext4}{10cm}

\vfill

\begin{blocgauche}
On ajoute un fragment, on redistribue, le répertoire ne change pas. C'est le cas qui devient 
le plus courant quand la collection grandit.
\end{blocgauche}

\end{frame}
%------------------------------------------------------------
\begin{frame}{À retenir\,: hachage extensible}

Le hachage extensible résout en partir le principal
défaut du hachage, l'absence de dynamicité.

\vfill

Le répertoire tend à croître de manière \alert{exponentielle}, ce qui peut soulever un problème
à terme.

\vfill
\begin{blocgauche}
Il reste une structure \alert{plaçante} qui doit être complétée par l'arbre B pour
des index secondaires.
\end{blocgauche}

\end{frame}
